\section{Rounding by pinching}

\IslamNote{Need to give proper credit to stackexchange and Overlapping qubits paper.}

\begin{fact}[Projector Icebreakers]\label{I-minus-P}
    If $P$ and $Q$ are (Hermitian) projectors, then:
    \begin{enumerate}
        \item $P^\bot = (I - P)$ is a (hermitian) projector;
        \item $\left(2P - I\right)$ is involutive (i.e. $\left(2P - I\right)^2 = I$);
        \item $P P^\bot = 0$;
        \item $[P, P^\bot] = 0$;
        \item $QPQ$ and $Q^\bot P Q^\bot$ are (hermitian) projectors; and
        \item $P' := QPQ + Q^\bot P Q^\bot$ are (hermitian) projectors.
    \end{enumerate}
Moreover, if $Q$ is a Hermitian projector, then:
\begin{enumerate}
\setcounter{enumi}{6}
    \item $U := (2Q - I)$ is Hermitian and unitary;
    \item $Q\ U = Q$.
    \item $Q^\bot\ U = -Q^\bot$.
    \item $\| Q\| = 1$.
\end{enumerate}
Furthermore, if $S$ and $T$ are two Hermitian Projectors orthogonal to each other (i.e. $\{S, T\}$ are orthonormal projectors), then:
\begin{enumerate}
\setcounter{enumi}{8}
    \item $S^\bot T = T S^\bot = T$; and
    \item $S^\bot T^\bot = T^\bot S^\bot = I - (S + T)$.
\end{enumerate}
Inductively, we also have for any $\{T_1, \dots, T_r\}$ set of orthonormal projectors, and any permutation $\pi: [r] \to [r]$
\begin{enumerate}
\setcounter{enumi}{10}
    \item $T_{\pi(1)}^\bot T_{\pi(2)}^\bot\dots T_{\pi(r)}^\bot = I - (\sum\limits_{k = 1}^r T_k)$.
\end{enumerate}


\end{fact}
\begin{proof}
    I will prove \# 6.
    \begin{align*}
        P'^2 &= \left(QPQ + Q^\bot P Q^\bot\right)^2\\
        &= (QPQ)^2 + QPQ Q^\bot P Q^\bot + Q^\bot P Q^\bot QPQ + \left(Q^\bot P Q^\bot\right)^2\\
        &= (QPQ)^2 + 0 + 0 + \left(Q^\bot P Q^\bot\right)^2 & QQ^\bot = Q^\bot Q = 0\\
        &= QPQ + Q^\bot P Q^\bot & QPQ \text{ and } Q^\bot P Q^\bot \text{ are projectors}\\
        &= P'.
    \end{align*}
\end{proof}

\begin{fact}[Spectral Theorem with a Cute Fact]
    Let $B$ be a Hermitian Operator $(B^\dagger = B)$ on a $d$-dimensional space. Then, there exist:
\begin{enumerate}
    \item a positive integer $0 < r \leq d$;
    \item a list of orthonormal projectors $\left(B_1, \dots, B_j, \dots, B_{r}\right)$; and 
    \item a list of non-zero real numbers (eigenvalues) in strictly increasing order of their absolute value $|b_1| < |b_2| <\dots < |b_r|$.
\end{enumerate}
 such that:
\begin{equation}
    B = \sum\limits_{j = 1}^{r} b_j B_j.
\end{equation}
Furthermore, define the ``support projection'' of $B$ to be $\Pi_B := \sum\limits_{j = 1}^r B_j$ and $\Pi_B^\bot := I - \Pi_B$. Then:
\begin{enumerate}
    \item $\Pi_B$ and $\Pi_B^\bot$ are hermitian projectors;
    \item $\Pi_B\ B = B\ \Pi_B = B$; and
    \item $\Pi_B^\bot\ B = B \ \Pi_B^\bot = 0$.
    \item $\|\Pi_B - B \| \leq \max_{j} | b_j - 1|\leq 1 + \| B \|$.
    \item If $b_1 > 0$, then $\|\Pi_B - B \| \leq \max_{j} |b_j - 1|\leq \| B \|$.
\end{enumerate}
\end{fact}
\begin{proof}
    \begin{align*}
        \|\Pi_B - B \| &= \|\Pi_B - \Pi_B\ B \| \\
        &= \|\Pi_B \left(I - B\right)\|\\
        &\leq \|\Pi_B\| \cdot \|I - B\|\\
        &= \|I - B\| & \|\Pi_B\| = 1\\
        &= \max_{j} |b_j - 1|\\
        &\leq 1 + \|B\|
    \end{align*}
\end{proof}



% \begin{proof}
%     \begin{align*}
%         (I - P)^2 &= I^2 - 2P + P^2\\
%         &= I - 2P + P & (P^2 = P)\\
%         &= (I - P).
%     \end{align*}
% \end{proof}

\begin{lemma}[A projector sandwich commutes with the bread.]
    Let $P, Q$ be projectors. Then,
    \[[QPQ, Q] = [QPQ, Q^\bot] = [Q^\bot P Q^\bot, Q^\bot] = [Q^\bot P Q^\bot, Q] = 0.\]
\end{lemma}
\begin{proof}
    \begin{align*}
        [QPQ, Q] &= QPQ^2 - Q^2PQ\\
        &= QPQ - QPQ\\
        &= 0.
    \end{align*}

\noindent To see why the second commutator is zero, observe:
\[
    [QPQ, Q^\bot] = QPQQ^\bot - Q^\bot QPQ = 0 - 0 = 0.
\]

\noindent The third and fourth commutators can be proven to be zero as a corollary by applying the above with the projector $Q^\bot$ in place of $Q$.
\end{proof}

\begin{fact}
    Let $A$ and $Q$ be $n \times n$ matrices where $Q$ is a projector (i.e. $Q^2 = Q$). Define:
    \begin{equation}
        X := QAQ + Q^\bot AQ^\bot
    \end{equation}
    
    \noindent Then, $[X, Q] = 0$.
    
    \noindent If $Q$ is a hermitian projector, then $\|X - A\| = \|[A,Q]\|$.
\end{fact}
\begin{proof}
    \begin{align*}
    \| X - A \| &= \| QAQ + Q^\bot A Q^\bot - A\|\\
    &= \| QAQ + (I-Q)A(I-Q) - A \|\\
    &= \| QAQ + A - AQ - QA + QAQ - A \|\\
    &= \| QAQ - AQ - QA + QAQ \|\\
    &= \|(Q - I)AQ - QA(I - Q) \|\\
    &= \|(I - Q)AQU + QA(I - Q) U \| & \cdot \times \|-U\| = 1\\
    &= \|(I - Q)AQ + QA(Q - I) \|\\
    &= \| AQ - QAQ + QAQ - QA\|\\
    &= \| AQ - QA\|\\
    &= \|[A, Q] \|.
\end{align*}
\end{proof}

\begin{theorem}[Pinching a Projector]
    Let $P$ and $Q$ be (hermitian) projectors and let 
    \[P' := QPQ + Q^\bot P Q^\bot.\]
    Then, $P'$ is a (hermitian) projector such that $[P', Q] = 0$. 
    
    \noindent Furthermore, if $P$ is any projector and $Q$ is a hermitian projector, then $\| P' - P \| = \|[P,Q] \| $.
\end{theorem}
\begin{proof}
$P'$ is a (hermitian) projector by Fact~\ref{I-minus-P}. We now prove it commutes with $Q$.

\begin{align*}
    [P', Q] &= [QPQ + Q^\bot P Q^\bot, Q]\\
    &= [QPQ, Q] + [Q^\bot P Q^\bot, Q]\\
    &= 0 + 0.
\end{align*}
When $Q$ is a hermitian projector, $U := (2Q - I)$ is unitary with $\|U\| = \|-U\| = 1$. %Further, note that $Q = \frac{I + U}{2}$ and $Q^\bot = \frac{I - U}{2}$.

\end{proof}

%\noindent Let's now look at what happens when you try to ``pinch'' using an oblique projector (bread).



\begin{lemma}[Pinching any matrix by a hermitian projector]
Let $A$ be an operator on a $d$-dimensional space. For a hermitian projector $Q$ on the same space, define:
\begin{equation}
    [A \overset{\text{pinch}}{\gets} Q] :=  Q A Q + (I-Q) A (I - Q).
\end{equation}

% \noindent According to the spectral theorem, there exist:
% \begin{enumerate}
%     \item a positive integer $0 <r \leq d$;
%     \item a list of orthonormal projectors $\left(A_1, \dots, A_j, \dots, A_{r}\right)$; and 
%     \item a list of real numbers (eigenvalues) in non-decreasing order of their absolute value $\left(a_1, \dots, a_j, \dots, a_{r}\right)$
% \end{enumerate}
%  such that:
% \begin{equation}
%     A = \sum\limits_{j = 1}^{r} a_j A_j
% \end{equation}

\noindent Then,
\begin{enumerate}
    \item $[A \overset{\text{pinch}}{\gets} Q]$ commutes with $Q$ i.e. $[[A \overset{\text{pinch}}{\gets} Q], Q] = 0$; and
    \item $\| [A \overset{\text{pinch}}{\gets} Q] - A\| = \|[A,Q]\|$.
\end{enumerate}

\noindent If $A$ is hermitian, then $[A \overset{\text{pinch}}{\gets} Q]$ is Hermitian.

\end{lemma}

\begin{proof}
To prove that $[A \overset{\text{pinch}}{\gets} Q]$ is Hermitian.

\begin{align*}
    [A \overset{\text{pinch}}{\gets} Q]^\dagger &=  \left(Q A Q + (I-Q) A (I - Q)\right)^\dagger \\
    &= Q^\dagger A^\dagger Q^\dagger + (I - Q)^\dagger A^\dagger (I-Q)^\dagger\\
    &= Q A Q + (I-Q) A (I - Q)\\
    &= [A \overset{\text{pinch}}{\gets} Q]
\end{align*}

\noindent The others follow from the previous results.

% \noindent Now, let's prove that $[A \overset{\text{pinch}}{\gets} Q]$ commutes with $Q$.

% \begin{align*}
%     \left[[A \overset{\text{pinch}}{\gets} Q], Q\right] &= \left[Q A Q + (I-Q) A (I - Q), Q\right]\\
%     &= \left[Q \left(\sum\limits_{j = 1}^{r} a_j A_j\right) Q + (I-Q) \left(\sum\limits_{j = 1}^{r} a_j A_j\right) (I - Q), Q\right]\\
%     &= \left[\sum\limits_{j = 1}^r a_j Q A_jQ + a_j(I - Q)A_j(I-Q), Q\right]\\
%     &= \sum\limits_{j = 1}^r a_j \left[Q A_jQ + (I - Q)A_j(I-Q), Q\right]\\
%     &= \sum\limits_{j = 1}^r a_j \left[[A_j \overset{\text{pinch}}{\gets} Q], Q\right]\\
%     &= \sum\limits_{j = 1}^r a_j \cdot 0 = 0.
% \end{align*}

% \noindent Now, we bound the distance.

% \begin{align*}
%     \| [A \overset{\text{pinch}}{\gets} Q] - A \| &=  \| Q A Q + (I-Q) A (I - Q) - A \|\\
%     &= \| Q \left(\sum\limits_{j = 1}^{r} a_j A_j\right) Q + (I-Q) \left(\sum\limits_{j = 1}^{r} a_j A_j\right) (I - Q) - \left(\sum\limits_{j = 1}^{r} a_j A_j\right) \|\\
%     &= \| \sum\limits_{j = 1}^{r} a_j \left(Q A_j Q + (I-Q) A_j (I - Q) - A_j\right) \|\\
%     &\leq \sum\limits_{j = 1}^{r} |a_j| \cdot \| \left(Q A_j Q + (I-Q) A_j (I - Q) - A_j\right) \|\\
%     &= \sum\limits_{j = 1}^{r} |a_j| \cdot  \| [A_j \overset{\text{pinch}}{\gets} Q] - A_j \|\\
%     &= \sum\limits_{j = 1}^{r} |a_j| \cdot  \| [A_j , Q] \|\\
%     &\leq r \cdot |a_r| \max_j{\| [A_j , Q] \|}\\
%     &\leq d \cdot \|A\|\|[A, Q]\| & \text{is that legal?}
% \end{align*}
    
\end{proof}




% Old stuff

\section{Old stuff}
\begin{lemma}[Pinching a hermitian operator by a set of orthonormal projectors]

Let $A$ be a hermitian operator with spectral decomposition:
\begin{equation}
    A = \sum\limits_{j = 1}^r a_j A_j
\end{equation}
and let $\mathcal{Q} = \left\{Q_1, \dots, Q_m\right\}$ be a set of orthonormal projectors.
Define:
\begin{equation}
    \Pi_{\mathcal{Q}} := \sum\limits_{Q \in \mathcal{Q}} Q \quad\quad\quad \text{ and } \quad\quad\quad \Pi_{\mathcal{Q}}^\bot := I - \Pi_{\mathcal{Q}}
\end{equation}
One can see that $\Pi_Q$ and $\Pi_Q^\bot$ are hermitian projectors. Then, define the pinching of $A$ by $\mathcal{Q}$ to be as follows:
\begin{align*}
    [A \overset{\text{pinch}}{\gets} \mathcal{Q}] &:= \left(\sum\limits_{j = 1}^m Q_j AQ_j\right) + \left(\prod\limits_{j = m}^1 Q_j^\bot\right) A  \left(\prod\limits_{j = 1}^m Q_j^\bot\right)\\
    &= \left(\sum\limits_{j = 1}^m Q_j AQ_j\right) + (I - Q_m) \dots (I-Q_1) A (I - Q_1) \dots (I - Q_m)\\
    &= \left(\sum\limits_{j = 1}^m Q_j AQ_j\right) + \Pi_{\mathcal{Q}}^\bot \ A\ \Pi_{\mathcal{Q}}^\bot
\end{align*}
Then, we can conclude the following:
\begin{enumerate}
    \item let $\pi: [m] \to[m]$ be any permutation, then:
    \begin{equation}
        [A \overset{\text{pinch}}{\gets} \mathcal{Q}] = \left[\left[\left[[A \overset{\text{pinch}}{\gets} Q_{\pi(1)}] \overset{\text{pinch}}{\gets} Q_{\pi(2)}\right]\dots\right] \overset{\text{pinch}}{\gets} Q_{\pi(m)}\right]
    \end{equation}
    \item $[A \overset{\text{pinch}}{\gets} \mathcal{Q}]$ is hermitian.
    \item For any $Q \in \mathcal{Q}$, $\left[[A \overset{\text{pinch}}{\gets} \mathcal{Q}],Q\right] = \left[[A \overset{\text{pinch}}{\gets} \mathcal{Q}],Q^\bot\right] = 0$
    \item $\left[[A \overset{\text{pinch}}{\gets} \mathcal{Q}],\Pi_{\mathcal{Q}}\right] = \left[[A \overset{\text{pinch}}{\gets} \mathcal{Q}],\Pi_{\mathcal{Q}}^\bot\right] = 0$.
    \item $\|[A \overset{\text{pinch}}{\gets} \mathcal{Q}] - A\| \leq $
\end{enumerate}
    
\end{lemma}
\begin{proof}
    The first one can be verified directly from the definition. To see why $[A \overset{\text{pinch}}{\gets} \mathcal{Q}]$ is hermitian, note:
    \begin{align*}
        [A \overset{\text{pinch}}{\gets} \mathcal{Q}]^\dagger &= \left(\left(\sum\limits_{j = 1}^m Q_j AQ_j\right) + \Pi_{\mathcal{Q}}^\bot \ A\ \Pi_{\mathcal{Q}}^\bot\right)^\dagger\\
        &= \left(\sum\limits_{j = 1}^m Q_j^\dagger A^\dagger Q_j^\dagger \right) + (\Pi_{\mathcal{Q}}^\bot)^\dagger \ A^\dagger\ (\Pi_{\mathcal{Q}}^\bot)^\dagger\\
        &= [A \overset{\text{pinch}}{\gets} \mathcal{Q}].
    \end{align*}
To see why $[A \overset{\text{pinch}}{\gets} \mathcal{Q}]$ commutes with any $Q_z \in \mathcal{Q}$, note the following.
\begin{align*}
    \left[[A \overset{\text{pinch}}{\gets} \mathcal{Q}], Q_z\right] &= \left(\sum\limits_{j = 1}^m Q_j AQ_j Q_z\right) + \Pi_{\mathcal{Q}}^\bot \ A\ \Pi_{\mathcal{Q}}^\bot Q_z - \left(\sum\limits_{j = 1}^m Q_z Q_j AQ_j\right) - Q_z\Pi_{\mathcal{Q}}^\bot \ A\ \Pi_{\mathcal{Q}}^\bot\\
    &= Q_z A Q_z - Q_z A Q_z = 0.
\end{align*}
Then $[A \overset{\text{pinch}}{\gets} \mathcal{Q}]$ commutes with $\Pi_{\mathcal{Q}}$ by linearity of the commutator.
Now, let's bound the distance. For ease of notation, let $X_{j}$ be $[A \overset{\text{pinch}}{\gets} \{Q_1, \dots, Q_j\}]$. Then, $X_m = [A \overset{\text{pinch}}{\gets} \mathcal{Q}]$.
\begin{align*}
    \|[A \overset{\text{pinch}}{\gets} \mathcal{Q}] - A\| &= \|[A \overset{\text{pinch}}{\gets} \mathcal{Q}] - X_{m-1} + X_{m-1} - A\|\\
    &\leq \|[X_{m-1} \overset{\text{pinch}}{\gets} Q_m] -X_{m-1}\| + \|X_{m-1} - A\| & \triangle  \text{ inequality}\\
    &\leq r |a_r| \max_{j}\| [A_j, Q_k]\| + \|X_{m-1} - A\|\\
    &\leq \sum\limits_{k = 1}^m r |a_r| \max_{j}\| [A_j, Q_k]\| & \triangle  \text{ inequality $m$ times}\\
    &\leq m \cdot r \cdot |a_r|\cdot \max_{j, k}\| [A_j, Q_k]\|\\
    &\leq m \cdot d \cdot \|A\|\cdot \max_{k}\| [A, Q_k]\|
\end{align*}
\end{proof}


\begin{lemma}[Pinching a hermitian operator by a hermitian operator]
Let $A = \sum\limits_{j = 1}^r a_j A_j$ and $B = \sum\limits_{k = 1}^m b_k B_k$ be the spectral decomposition of two Hermitian operators $A$ and $B$ on a $d$-dimensional space with bounded norm $\|A\| \leq a$ and $\|B\| \leq b$. Define $\mathcal{B} := \{B_k\}$ and define:
\begin{equation}
    [A \overset{\text{pinch}}{\gets} B] :=  [A \overset{\text{pinch}}{\gets} \mathcal{B}]
\end{equation}

\noindent Then,
\begin{enumerate}
    \item $[A \overset{\text{pinch}}{\gets} B]$ is Hermitian;
    \item $[A \overset{\text{pinch}}{\gets} B]$ commutes with $B$ i.e. $[[A \overset{\text{pinch}}{\gets} B], B] = 0$; and
    \item $\| A - [A \overset{\text{pinch}}{\gets} B] \| \leq$ (same bound from previous lemma).
\end{enumerate}
\end{lemma}

\begin{proof}

\begin{align*}
     [[A \overset{\text{pinch}}{\gets} B], B] &= \left[\sum\limits_{k} B_k A B_k + \Pi_{\mathcal{Q}}^\bot \ A\ \Pi_{\mathcal{Q}}^\bot + , B\right]\\
     &= \sum\limits_{k, q} b_q\left[B_k A B_k + \Pi_{\mathcal{Q}}^\bot \ A\ \Pi_{\mathcal{Q}}^\bot + , B_q\right] = 0.
    % [\sum\limits_{k} B_k A B_k, B] + \left[\Pi_{\mathcal{Q}}^\bot \ A\ \Pi_{\mathcal{Q}}^\bot, B\right]\\
    % &= \sum\limits_{k, q} [B_k A B_k, b_q B_q]\\
    % &= \sum\limits_{k, q} b_q \delta_{k,q} [B_k A B_k, B_k]\\
    % &= 0.
\end{align*}

The distance is the same one from the previous lemma.

\end{proof}

\begin{example}[Oblique Projector]
    Consider the projector:
    \begin{equation}
        Q = E_{11} + E_{12} = \begin{pmatrix}
            1 & 1\\
            0 & 0
        \end{pmatrix}.
    \end{equation}
    One can verify that it is a projector because $Q^2 = Q$. However, it is \textbf{oblique} and not orthogonal (Hermitian) because:
    \begin{equation}
        Q^\dagger = \left(E_{11} + E_{12}\right)^\dagger = E_{11} + E_{21} = \begin{pmatrix}
            1 & 1\\
            0 & 0
        \end{pmatrix}^\dagger = \begin{pmatrix}
            1 & 0\\
            1 & 0
        \end{pmatrix} \neq Q.
    \end{equation}
    By looking at $(2Q - I)$ and $(2Q - I)^\dagger$, we will see that neither of them is unitary nor Hermitian.
    \begin{equation}
        2Q - I = 2 \cdot\begin{pmatrix}
            1 & 1\\
            0 & 0
        \end{pmatrix} - \begin{pmatrix}
            1 & 0\\
            0 & 1
        \end{pmatrix} = \begin{pmatrix}
            1 & 2\\
            0 & -1
        \end{pmatrix}.
    \end{equation}
    \begin{equation}
        (2Q-I)^\dagger = \begin{pmatrix}
            1 & 0\\
            2 & -1
        \end{pmatrix} \neq \left(2Q - I\right).
    \end{equation}
\end{example}

\begin{example}[Distance blowing larger than the commutator when $Q$ is not Hermitian; could it still be bounded within a factor???; need to check that]
    Consider the following two oblique non-Hermitian projectors:
    \begin{equation}
        P = E_{11} + E_{21} = \begin{pmatrix}
            1 & 0 \\1 & 0
        \end{pmatrix}\quad\quad\quad Q = E_{11} + E_{12} = \begin{pmatrix}
            1 & 1\\0 &0
        \end{pmatrix}
    \end{equation}
    \begin{equation}
        PQ = E_{11} + E_{12} + E_{21} + E_{22} = \begin{pmatrix}
            1 & 1\\
            1 & 1
        \end{pmatrix}
    \end{equation}
    \begin{equation}
        QP = 2E_{11} = \begin{pmatrix}
            2 & 0\\0 & 0
        \end{pmatrix}
    \end{equation}
    \begin{equation}
        [P,Q] = PQ - QP = \begin{pmatrix}
            -1 & 1\\ 1 & 1
        \end{pmatrix}
    \end{equation}
Let's now compute the operator norm for this matrix.
\begin{align*}
    [P,Q]^\dagger[P,Q] = [P,Q]^2= 2I
\end{align*}
\noindent Thus, $2$ is the largest eigenvalue of $[P,Q]^\dagger[P,Q]$ and thus $\|[P,Q]\| = \sqrt{2}$.

\noindent Now, let's compute the operator norm distance between $P'$ and $P$.
\begin{align*}
    D = P' - P &=  QPQ - (I-Q)P(I-Q) - P \\
    &= 2QPQ - PQ - QP \\
    &= 2QPQ - (PQ + QP)\\
    &= 2\cdot 2E_{11} \left(E_{11} + E_{12}\right) - \begin{pmatrix}
        3 & 1\\1&1
    \end{pmatrix}\\
    &= 4\left(E_{11} + E_{12} \right) - \begin{pmatrix}
        3 & 1\\1&1
    \end{pmatrix} \\
    &= \begin{pmatrix}
        4 & 4\\0 & 0
    \end{pmatrix} - \begin{pmatrix}
        3 & 1\\ 1 & 1
    \end{pmatrix} \\
    &= \begin{pmatrix}
        1 & 3\\ -1 & -1
    \end{pmatrix} .
\end{align*}
\begin{align*}
    D^\dagger D &= \begin{pmatrix}
        1 & -1 \\ 3 & -1
    \end{pmatrix} \begin{pmatrix}
        1 & 3\\-1 & -1
    \end{pmatrix} = \begin{pmatrix}2 & 4\\4&10
    \end{pmatrix}
\end{align*}
The eigenvalues of $D^\dagger D$ are $6 \pm 4\sqrt{2}$. Therefore, the operator norm of $D = P' - P$ is $\sqrt{6 + 4\sqrt{2}} > 3.41 > \sqrt{2}$.
\end{example}

\begin{lemma}
    Suppose $A, B$ are 
\end{lemma}



\newpage
\appendix
\section{OLD: Factorizing algebras}

\newcommand{\calA}{\mathcal{A}}
\newcommand{\calL}{\mathcal{L}}

Here all algebras are unital star algebras. 

\begin{definition}
    An algebra $\mathcal{A}$ is ``algebraically irreducible'' if its center consists of only the multiples of identity.
\end{definition}
\begin{definition}
    An algebra $\mathcal{A} \subseteq \mathcal{L}(V)$ for some vector space $V$ is irreducible if it has no invariant subspaces.
\end{definition}
\begin{claim}
If $\calA \subseteq \calL(V)$ is irreducible, then it is algebraically irreducible.
\end{claim}
\begin{proof}
    Suppose not, i.e. suppose that $\calA$ is not algebraically irreducible. Then there is some element $z \in Z(\calA)$ that is not equal to a multiple of identity. This $z$ has an eigenvalue $\lambda$ and $z - \lambda I \in \calA$ and non-invertible. Consider the kernel of $z - \lambda I$. This is an invariant subspace of $\calA$: suppose $v \in \ker(z -\lambda I)$. Then for all $a \in \calA$,
    \[ (z - \lambda I) a v = a (z - \lambda I)v = 0,\]
    so $av \in \ker(z - \lambda I)$ as well. Moreover, since $z$ was not itself a multiple of identity, this invariant subspace is not trivial, which contradicts the irreduciblity of $\calA$.
\end{proof}

\begin{theorem}[Burnside]
For a finite dimensional complex vector space $V$, the only irreducible subalgebra of $\calL(V)$ is $\calL(V)$ itself.
\end{theorem}

\begin{lemma}[Fact 2.13 of Aharonov and Eldar]
    Let $\calA \subseteq \calL(V)$ be an algebra. Then there is a decomposition of $V$ into orthogonal subspaces $V_\alpha$, where each $V_\alpha$ is a tensor product of two vector spaces $V_\alpha = V^1_\alpha \otimes V^2_\alpha$ such that
    \[ \calA \cong \bigoplus_{\alpha} \calL(V^1_\alpha) \otimes \mathbf{I}(V^2_\alpha). \]
    
\end{lemma}
\begin{proof}
    Decompose $V$ into invariant subspaces for $\calA$ as follows: let $V_0$ be an invariant subspace of $V$ under $\calA$ of minimal dimension that is not equal to $\{0\}$ (since $V$ itself is an invariant subspace, such a $V_0$ is guaranteed to exist). Write $V = V_0 \oplus V_0^\perp$. Then $V_0^\perp$ is an invariant subspace for $\calA$ as well. Moreover, since $V_0$ was of minimal dimension, $\calA$ must be irreducible restricted to $V_0$. Inductively continue this procedure on $V_0^\perp$ until there's nothing left. This gives us a decomposition
    \[ V = V_0 \oplus V_1 \oplus \dots \]
    where the $V_i$s are all orthogonal, and $\calA$ is irreducible on each $V_i$. By Burnside, this means that for each $V_i$, $\calA|_{V_i} = \calL(V_i)$. 

    Let us now group together subspaces $V_i$ according to isomporphism as $\calA$-modules: $V_i \cong V_j$ as an $\calA$ module if there exists an ``intertwiner map'' $T_{ij}: V_i \to V_j $ which is a bijective linear map such that 
    \[ \forall v \in V_i, a \in \calA, T_{ij}(a v) = a T_{ij}(v).\] Let us label the equivalence classes under this relation by $\alpha$, and the multiplicity of each class by $m_\alpha$. Then we can write
    \[ V = \bigoplus_\alpha V^1_\alpha \otimes V^2_\alpha, \]
    where $V^1_\alpha$ is a representative of this isomorphism class and $V^2_\alpha \cong \mathbb{C}^{m_\alpha}$. Now we want to claim that
    \[ \calA \cong \bigoplus_\alpha \calL(V^1_\alpha) \otimes \mathbf{I}(V^2_\alpha).\]
    CHAT GPT SAYS try Schur's lemma here. 

    % Let's now look at the center $\calA$, which we denote $Z(\calA)$. It is not hard to see that if $M \in Z(\calA)$, then
    % \[ M = \bigoplus_i \gamma_i \Pi_i,\]
    % where $\Pi_i$ is the projector onto $V_i$. This is because anything that doesn't look like a multiple of identity when restricted to one of the $V_i$s will fail to commute with some element of $\calA$ (since we know that $\calA |_{V_i} = \calL(V_i)$). Moreover, every element of $\calA$ that has this form is in $Z(\calA)$. 



    % \textcolor{red}{Now, let us group together $V_is$ that have the same dimension. For each unique dimension $d_\alpha$, let $m_\alpha$ be its multiplicity, and let $V_\alpha$ be the direct sum of all $V_i$s with dimension $d_\alpha$. Then we have that
    % \[ V_\alpha = \underbrace{V^1_\alpha}_{\cong \mathbb{C}^{d_\alpha}} \otimes \underbrace{V^2_\alpha}_{\cong \mathbb{C}^{m_\alpha}} \]
    % Moreover we get that
    % \[ \calA \cong \bigoplus_\alpha \calL(V^1_\alpha) \otimes \mathbb{I}(V^2_\alpha). \]}
\end{proof}

\pagebreak


\begin{claim}
    Let $\calA \subseteq \calL(V)$ be algebraically irreducible and let $T \subseteq V $ be an invariant subspace for $\calA$ such that $a_1, a_2$ are equal in their action on $T$. Then $a_1 = a_2$.
\end{claim}
\begin{proof}
    Suppose $a_1 \neq a_2$. Then $a_1 - a_2$ is nonzero but its kernel contains all of $T$. In particular, $(a_1 - a_2)$ commutes with $\calA$ on $T$. 
\end{proof}

\begin{claim}
    Let $\calA \subseteq \calL(V)$ be an algebra over a vector space $V$ and let $S$ be an invariant subspace of $\calA$. LALALALAAL
\end{claim}

\begin{lemma}
    Let $\mathcal{A}$ be an ``algebraically'' irreducible subalgebra of the algebra $\mathcal{L}(V)$ of linear operators on a finite-dimensional complex vector space $V$. Then we can decompose $V$ and $\mathcal{A}$ as
    \begin{align} 
    V &\approx S \otimes \mathbb{C}^k \\
    \mathcal{A} &\approx \mathcal{L}(S) \otimes I_k.
    \end{align}
\end{lemma}
This is Fact~2.12 of Aharonov-Eldar.
\begin{proof}
    Pick an arbitrary vector $v_1 \in V$. Let $S_1 = \mathcal{A}(v_1)$, i.e. the subspace consisting of all vectors that can be written as $a v_1$ for some $a \in \calA$. 

    SDFLSKJDFLSDKJFLSDKFJ

    We have decomposed 
    \[ V = \bigoplus_i S_i,\]
    where each $S_i$ is irreducible. Let's group them together into copies of the same space:
    \[V = \bigoplus_k S_k \otimes I_{n_k}, \]
    where $n_k$ counts the multiplicty of the space $S_k$.
    Now because of irreducibility, we have a homomorphism from $\calA$ to $\calL(S_k)$ for each $S_k$. 
\end{proof}
\newpage

